\documentclass[11pt]{article}
\usepackage[margin=1in]{geometry}
\usepackage{booktabs}
\usepackage{listings}
\usepackage{xcolor}
\usepackage{hyperref}
\usepackage{enumitem}
\usepackage{array}

\lstset{
  basicstyle=\ttfamily\small,
  frame=single,
  breaklines=true,
  columns=fullflexible,
  backgroundcolor=\color{white},
  commentstyle=\color{gray},
  showstringspaces=false,
  xleftmargin=2pt,
  xrightmargin=2pt
}

\title{\textbf{RustFS: Project Report}\\
\large A Journaled, Crash-Safe ext2-Style Filesystem in Rust,\\
Verified via a Real FUSE Mount and Real \texttt{fsck} Runs}
\author{Hemal Kisku}
\date{}

\begin{document}
\maketitle

\section{System Description}

\subsection{What we implemented}

We implemented a FUSE-based filesystem, \textbf{rustfs}, built from scratch in Rust and modeled
after the classic ext2 on-disk layout. The system is not a toy simulation of a filesystem's data
structures — it is a real block device image, mountable with the Linux kernel's FUSE interface,
that a shell, \texttt{dd}, \texttt{ls}, or any other unmodified POSIX program can read from and
write to.

The project consists of four crates sharing one core library:
\begin{itemize}
    \item \textbf{libfs} — the on-disk layout (superblock, inode, directory entry, journal header
    structs), a block-device abstraction with an LRU write-back cache, a bit-level allocator for
    inodes and data blocks, an inode table reader/writer, a directory-entry store, and a
    write-ahead journal.
    \item \textbf{mkfs} — formats a blank image with a superblock, both bitmaps, a zeroed inode
    table, and a root directory.
    \item \textbf{rustfs} (the FUSE daemon) — mounts an image and implements 16 POSIX operations
    (\texttt{lookup}, \texttt{getattr}, \texttt{setattr}, \texttt{readdir}, \texttt{create},
    \texttt{read}, \texttt{write}, \texttt{mkdir}, \texttt{rmdir}, \texttt{unlink},
    \texttt{rename}, \texttt{symlink}, \texttt{readlink}, \texttt{link}, \texttt{destroy}, plus
    initialization) against that image.
    \item \textbf{fsck} — an independent, from-scratch consistency checker that walks the same
    on-disk image without going through the FUSE daemon at all, and reports (rather than silently
    trusts) whether the metadata is internally consistent.
\end{itemize}

\subsection{Why it is useful}

Implementing a filesystem's data structures is the easy 80\% — inode tables and bitmaps are
straightforward once the layout is fixed. The parts that actually determine whether a filesystem
is trustworthy are the parts that are hard to get right by inspection alone:

\begin{itemize}
    \item \textbf{Crash consistency.} A metadata update that touches several blocks (allocate an
    inode, update a bitmap, link a directory entry) must appear atomic even if the process is
    killed mid-update. This cannot be verified by reading the code — it has to be verified by
    actually killing the process mid-update and checking what is left behind.
    \item \textbf{Consistency, not just correctness of the happy path.} A filesystem can pass
    every functional test and still leak an inode, double-allocate a block, or corrupt a bitmap
    counter under conditions the happy-path tests never exercise. \texttt{fsck} exists precisely
    to catch that class of bug independently of whatever the FUSE daemon believes about its own
    state.
    \item \textbf{Real measurement over assumed complexity.} A cost claim such as ``directory
    lookup is $O(n)$'' or ``fsck scales with used inodes'' is an assumption until it is measured.
    Every performance claim in this report is a number produced by actually running the code
    against a real mounted filesystem, not an estimate of the data structure's asymptotic shape.
\end{itemize}

This project was built to internalize those three properties — journaled crash safety, an
independent consistency checker, and measurement-verified performance claims — the same way a
production filesystem is expected to justify them, at a scale small enough to implement, break,
and fix by hand.

\section{Analysis Used}

\subsection{Dimensions of precision}

Several places in the design deliberately trade generality for a bounded, predictable cost, mirroring
the same kind of scope discipline production filesystems apply at much larger scale:

\begin{itemize}
    \item \textbf{Bounded block-pointer scheme.} Each inode has 12 direct pointers, one single-indirect
    pointer, and one double-indirect pointer — no triple indirect. At a 4\,KiB block size this caps a
    file at $12 + 512 + 512^2 \approx 262{,}656$ blocks ($\approx$1.03\,GiB), which is far larger than
    anything the 100\,MiB seed image can hold. Triple indirect was deliberately not implemented: it would
    add a third layer of pointer-chasing code for a ceiling the target image size can never reach.
    \item \textbf{Fixed-size, bounded journal.} The journal supports at most \texttt{MAX\_JOURNAL\_BLOCKS = 4}
    data blocks per transaction, written to a fixed 8-block on-disk region (one header block plus four data
    blocks) rather than a general-purpose variable-length circular log. This is enough to atomically cover
    every metadata operation the daemon performs (bitmap + superblock + inode-table + one directory block),
    and a fixed region means recovery only ever has to check one header, not scan for log boundaries.
    \item \textbf{Fixed-capacity LRU cache.} \texttt{Disk} caches exactly 256 blocks (1\,MiB) with
    write-back semantics rather than a dynamically-sized cache. Eviction flushes dirty entries before
    dropping them, so capacity is a memory/hit-rate tradeoff, never a correctness one.
    \item \textbf{Linked-list directory format.} Directory entries are a singly-linked list of
    variable-length records packed into one 4\,KiB block (\texttt{inode.direct[0]}, no external growth) — the
    same layout ext2 itself uses before hashed directories. This is $O(n)$ per lookup by construction; Section
    4.2 measures that cost directly rather than assuming it.
    \item \textbf{\texttt{fsck} scans by capacity, not by usage.} The inode-bitmap check iterates
    \texttt{0..total\_inodes} (a fixed 2048-slot range) regardless of how many inodes are actually allocated.
    This was a deliberate simplicity choice — it avoids needing a second, separately-maintained count of live
    inodes — at the cost of \texttt{fsck}'s runtime being bound by total filesystem capacity rather than by how
    full the filesystem is. Section 4.2 measures this tradeoff explicitly.
\end{itemize}

\subsection{Assumptions}

\begin{itemize}
    \item \textbf{Single physical writer.} All disk access goes through one \texttt{Arc<Mutex<Disk>>}
    shared by the FUSE daemon's callback handlers. This serializes every block read/write across
    concurrent FUSE requests — correct and simple, at the cost of no parallel I/O within one mount.
    \item \textbf{Fixed on-disk geometry.} The layout (2048 inodes across 128 inode-table blocks, one
    global 8-block journal region, a 100\,MiB seed image) is chosen at \texttt{mkfs} time and is not
    resizable afterward. There is no online or offline resize tool.
    \item \textbf{At most one pending transaction.} Journal recovery assumes a single global journal
    region can hold at most one committed-but-unapplied transaction at a time — correct because
    \texttt{begin\_transaction} is only ever called from within one already-serialized FUSE handler, but
    it would not extend to a design that pipelines multiple in-flight transactions.
    \item \textbf{Flat, C-compatible on-disk structs.} Every on-disk structure (\texttt{Superblock},
    \texttt{Inode}, \texttt{DirEntryHeader}, \texttt{JournalHeader}) is \texttt{\#[repr(C)]} and
    (de)serialized with raw \texttt{std::ptr::read}/\texttt{std::ptr::write} rather than a serialization
    crate — a zero-copy design that assumes the reader and writer always agree on the exact struct layout
    (no versioning, no forward/backward compatibility).
\end{itemize}

\section{Implementation Details}

\subsection{Write path: cache, journal, then FUSE}

A metadata-mutating FUSE call (e.g. \texttt{create}) follows the same shape every time: compute the
new in-memory state (allocated inode, updated bitmap, updated superblock counters, updated directory
block), hand every touched block to \texttt{journal.begin\_transaction(\&[\ldots])} \emph{before} any
of those blocks are written to their real locations, then apply the real writes, then call
\texttt{journal.commit\_complete()} to clear the journal. On mount, \texttt{Journal::recover()} checks
the journal header for a committed-but-not-yet-cleared transaction and replays it before the daemon
serves its first request — so a crash between the journal write and the real writes is always
recoverable, and a crash before the journal write simply never happened as far as the filesystem is
concerned.

\subsection{Before/After examples}

\textbf{Snippet 1 --- Crash injected mid-transaction, recovered on remount.} A \texttt{CRASH\_TEST}
environment variable was added at the exact point between the journal commit and the real writes, to
make the crash window reproducible instead of relying on timing an external kill signal:

\begin{lstlisting}
journal.begin_transaction(&[
    (inode_bitmap_block, new_inode_bitmap),
    (0, sb_block),
    (inode_block_idx, inode_table_block_data),
    (dir_block_num, dir_block_data),
]);

if std::env::var("CRASH_TEST").is_ok() {
    std::process::exit(1);   // simulate a crash after commit, before real writes
}
// ... real writes to bitmap / superblock / inode table / directory ...
journal.commit_complete();
\end{lstlisting}
Running \texttt{CRASH\_TEST=1 cargo run --bin rustfs -- myfs.img /tmp/mnt} against a \texttt{create}
kills the daemon after the journal is committed but before any real block is written. Remounting
without \texttt{CRASH\_TEST} set triggers \texttt{Journal::recover()}, which replays all four blocks
from the journal and leaves the image identical to a run that never crashed.

\textbf{Snippet 2 --- Indirect addressing before and after the fix.} Before
\texttt{d093a25}, writes only ever touched \texttt{inode.direct[0..12]} — files were silently capped
at 48\,KiB with no error. \texttt{get\_or\_alloc\_block} now dispatches on the logical block index:
\begin{lstlisting}
fn get_or_alloc_block(&mut self, inode: &mut Inode, block_index: usize) -> Option<u64> {
    if block_index < 12 {
        // direct
    } else if block_index < 12 + 512 {
        // single indirect: allocate/zero the pointer block on first use, then index into it
    } else if block_index < 12 + 512 + 512 * 512 {
        // double indirect: allocate outer block, then inner block, then the data block
    } else {
        None
    }
}
\end{lstlisting}
A 20\,MiB \texttt{dd} write (Section 4.2, Benchmark 1) exercises the single-indirect path directly and
would have failed outright under the pre-fix direct-only code.

\textbf{Snippet 3 --- Duplicate-block detection, before and after.} \texttt{fsck} originally recorded
referenced blocks in a \texttt{HashSet<u64>}:
\begin{lstlisting}
// before: two inodes pointing at the same block collapse into one set entry —
// the second owner is invisible, so duplicate allocation can never be reported.
if block != 0 { referenced_blocks.insert(block); }
\end{lstlisting}
\begin{lstlisting}
// after: every owner is recorded, so a block with more than one owner is detectable,
// and re-walking the same inode twice (e.g. via two directory paths) is guarded separately
// so it isn't double-counted as its own "duplicate".
block_owner.entry(block).or_default().push(inode_num);
\end{lstlisting}
This also folded in bounds checking (a block index outside \texttt{[data\_start\_block, total\_blocks)}
is reported immediately instead of being followed into an out-of-range read), and a
\texttt{collected\_inodes} guard so an inode reachable by more than one directory path is only
scanned for its own block pointers once.

\subsection{Key challenges}

\begin{enumerate}
    \item \textbf{Journal did not cover directory writes.} The journal originally wrapped the inode
    bitmap, superblock, and inode-table writes for \texttt{create}, but the directory-entry write that
    actually links the new inode into its parent directory happened afterward, unjournaled. A crash in
    that window left a fully-allocated inode with no directory entry pointing at it — an orphaned inode,
    invisible to \texttt{ls} but consuming an inode slot forever. Fixed in \texttt{e66f46d} by adding
    \texttt{DirStore::prepare\_add\_entry}, which computes the modified directory block \emph{without}
    writing it, so it can be folded into the same \texttt{begin\_transaction} call as the other three
    blocks instead of being written separately afterward. This is also precisely the failure mode the
    \texttt{fsck} orphaned-inode check (added next, in \texttt{ae5872f}) was built to detect.
    \item \textbf{\texttt{fsck} traversed the same inode and the same block more than once.} Two related
    bugs, fixed together in \texttt{a25af53}: (a) a directory reachable by more than one path (or visited
    while walking both its own entries and a hard link) had its block pointers collected once per visit,
    which double-counted legitimately-shared blocks as ``duplicates''; (b) the referenced-block set was a
    plain \texttt{HashSet}, which cannot distinguish ``referenced once'' from ``referenced by two different
    inodes'' — the actual definition of a duplicate-allocation bug. Fixed by tracking a
    \texttt{HashMap<block, Vec<owning\_inode>>} instead of a set, and gating block collection per inode
    behind a \texttt{collected\_inodes} guard (Snippet 3 above).
    \item \textbf{Single- and double-indirect blocks were never wired into the write path.} Only direct
    pointers were populated by \texttt{write()}/\texttt{read()}, silently truncating any file past
    48\,KiB with no reported error — the worst kind of bug, since it fails quietly rather than loudly.
    Fixed in \texttt{d093a25} by implementing \texttt{get\_or\_alloc\_block} and
    \texttt{get\_block\_readonly} with the three-tier dispatch shown in Snippet 2.
    \item \textbf{\texttt{fsck}'s diagnostic dump printed struct offsets instead of field values.}
    \texttt{std::mem::offset\_of!(Superblock, field)} returns the field's \emph{byte offset within the
    struct}, not its value. This was accidentally used in the startup diagnostic dump, so a real
    2048-inode filesystem printed \texttt{total\_inodes 24} (an offset) instead of \texttt{2048} (the
    actual count) — a purely cosmetic bug, since every actual consistency check downstream read the
    correct field values via ordinary struct access; only the printed dump was wrong. Caught while
    diagnosing the \texttt{fsck} scan-time benchmark (Section 4.2, Benchmark 4) and fixed by replacing
    every \texttt{offset\_of!} call with the field itself.
    \item \textbf{Benchmark methodology artifacts had to be ruled out before trusting a number,
    for every one of the four benchmarks.} An early metadata-throughput pass showed delete running
    17$\times$ slower than create (373 vs.\ 5{,}997 ops/s). The filesystem was not the cause: the delete
    loop invoked \texttt{rm} once per file (forking a new process every iteration) while the create loop
    used bash's \texttt{: >} builtin (no fork). Running the identical \texttt{rm} loop against the host's
    own ext4 reproduced the same $\sim$0.4--0.8\,s overhead with no filesystem code involved at all,
    confirming the gap was shell process-startup cost, not rustfs. Corrected by measuring deletes with
    \texttt{find -delete} (a single process), which brought both numbers into the same range
    ($\sim$5{,}800--6{,}000 ops/s). The sequential-I/O benchmark required a matching discipline: a fresh
    image per run, remounted between write and read, specifically to drop both the in-process block
    cache and the kernel page cache so a warm-cache read was never mistaken for real disk throughput.
\end{enumerate}

\section{Result Tables}

\subsection{Functional correctness --- what \texttt{fsck} checks}

\texttt{fsck} is a from-scratch, independent walk of the on-disk image (it does not go through the
FUSE daemon or its cache), and reports every inconsistency it finds rather than stopping at the first
one:

\begin{center}
\begin{tabular}{p{5.4cm}p{8.4cm}}
\toprule
\textbf{Check} & \textbf{What it catches} \\
\midrule
Magic number & Corrupt or wrong image \\
Block pointer validity & Pointers outside the data region \\
Duplicate block allocation & Two inodes sharing a data block \\
\texttt{.} / \texttt{..} correctness & Directory self/parent links wrong \\
Directory cycle detection & Infinite loops in the directory tree \\
Inode bitmap consistency & Inodes referenced but marked free; orphaned inodes \\
Data bitmap consistency & Leaked blocks; blocks referenced but marked free \\
Hard link counts & \texttt{hard\_links} field doesn't match actual directory references \\
Superblock counters & \texttt{free\_blocks} / \texttt{free\_inodes} out of sync with bitmaps \\
\bottomrule
\end{tabular}
\end{center}

\subsection{Benchmarks}

Four benchmarks, 5 runs each, against a 100\,MiB image on a local FUSE mount. Every number below
required ruling out a measurement artifact first (Section 3.3, item 5); the methodology described
there is what makes these numbers trustworthy rather than just fast-looking.

\textbf{Benchmark 1 --- Sequential I/O throughput.} 20\,MiB file, \texttt{dd}, fresh image per run,
remounted between write and read.

\begin{center}
\begin{tabular}{lrr}
\toprule
 & \textbf{Mean} & \textbf{Stddev} \\
\midrule
Write & 74.28 MiB/s & 5.33 \\
Read  & 414.14 MiB/s & 100.5 \\
\bottomrule
\end{tabular}
\end{center}

\textbf{Benchmark 2 --- Metadata operation throughput.} 200 file creates/deletes per run through the
FUSE mount, measured with the corrected (single-process) methodology from Section 3.3, item 5.

\begin{center}
\begin{tabular}{lr}
\toprule
 & \textbf{Mean} \\
\midrule
Create & $\sim$5{,}997 ops/s \\
Delete & $\sim$5{,}826 ops/s \\
\bottomrule
\end{tabular}
\end{center}

\textbf{Benchmark 3 --- Directory lookup latency vs.\ entry count.} In-process microbenchmark
(\texttt{libfs/examples/bench\_dir\_lookup.rs}), bypassing FUSE entirely to isolate
\texttt{DirStore::look\_up()} from kernel dispatch overhead. 100{,}000 iterations per data point, 5
runs averaged, entry counts from 10 up to $\sim$300 (the block's practical capacity):
\[
\text{latency(ns)} \approx 207 + 20.7 \times \text{entry\_count} \qquad (R^2 > 0.999)
\]
This confirms the $O(n)$ cost predicted by the linked-list directory format (Section 2.1) directly,
rather than assuming it from the data structure alone, and marks hashed or indexed directory entries
as the natural next step if larger directories are ever needed.

\textbf{Benchmark 4 --- \texttt{fsck} scan time vs.\ file count.} \texttt{fsck} timed against images
pre-populated with 0, 50, 100, 150, and 200 files, 5 runs each.

\begin{center}
\begin{tabular}{lr}
\toprule
\textbf{Files} & \textbf{Mean} \\
\midrule
0   & 3.4 ms \\
50  & 3.0 ms \\
100 & 2.4 ms \\
150 & 2.4 ms \\
200 & 2.8 ms \\
\bottomrule
\end{tabular}
\end{center}
Flat, not linear — confirming the Section 2.1 prediction that \texttt{fsck}'s cost is bound by total
filesystem capacity (a fixed 2048-inode scan range), not by how full the filesystem is.

\subsection{Real bugs found and their measured/observed effect}

\begin{center}
\begin{tabular}{p{3.6cm}p{4.6cm}p{4.6cm}}
\toprule
\textbf{Bug} & \textbf{Symptom before fix} & \textbf{Result after fix} \\
\midrule
Journal excluded directory writes & Crash between inode alloc and dir-link produced an orphaned inode & Directory write folded into the same journaled transaction; recoverable across the crash window (Snippet 1) \\
\texttt{fsck} duplicate-block tracking used a \texttt{HashSet} & Two inodes sharing a block, or one inode walked twice, were indistinguishable & \texttt{HashMap<block,\allowbreak Vec<owner>>} + per-inode collection guard; duplicates now attributable to actual owners \\
Indirect/double-indirect never wired into writes & Files silently truncated at 48\,KiB with no error & Three-tier \texttt{get\_or\_alloc\_block} dispatch; 20\,MiB sequential write (Benchmark 1) exercises the fix directly \\
\texttt{fsck} diagnostic dump used \texttt{offset\_of!} & Startup dump printed byte offsets (e.g.\ \texttt{total\_inodes 24}) instead of real values & Dump prints actual field values; consistency-check logic itself was never affected \\
Delete-benchmark artifact & Delete measured 17$\times$ slower than create (373 vs.\ 5{,}997 ops/s) & Traced to per-file \texttt{rm} fork cost (reproduced on host ext4); \texttt{find -delete} brings both to $\sim$5{,}800--6{,}000 ops/s \\
\bottomrule
\end{tabular}
\end{center}

\subsection{Discussion of results}

\begin{enumerate}
    \item \textbf{Read throughput exceeds write by $>5\times$ mainly because of the write-back cache
    and the 20\,MiB size relative to the 1\,MiB cache}: writes pay real disk I/O as dirty blocks are
    evicted, while a chunk of the subsequent read is served from blocks already resident. This is
    reported as-measured rather than attributed without evidence — a full page-cache-bypassing
    re-run per read would be needed to separate ``fast because cached'' from ``fast because the read
    path is simply cheaper,'' and is noted here as the natural follow-up rather than assumed.
    \item \textbf{The $O(n)$ directory-lookup cost is a direct, measured consequence of the format
    choice in Section 2.1}, not an incidental slowdown — the same linked-list layout that keeps
    \texttt{add\_entry}/\texttt{remove\_entry} simple is what produces the linear slope. The
    benchmark exists specifically to make that tradeoff visible instead of implicit.
    \item \textbf{\texttt{fsck}'s flat scan time is the direct, measured consequence of the capacity-bound
    design choice in Section 2.1}: iterating a fixed 2048-slot range costs the same whether 0 or 200
    files are populated, which is expected given the code, but was verified rather than assumed.
    \item \textbf{Every bug in Section 4.3 was found by cross-checking two independent views of the
    same on-disk state} — the FUSE daemon's writes against \texttt{fsck}'s independent walk, or a
    benchmark's number against a second, differently-instrumented measurement of the same operation —
    rather than by code inspection alone. This is the same principle the crash-injection harness
    (\texttt{CRASH\_TEST}, Snippet 1) applies to durability: don't trust that a recovery path works,
    trigger the crash and check what's actually on disk afterward.
\end{enumerate}

\vspace{1em}
\noindent\rule{\textwidth}{0.4pt}\\
\footnotesize
Stack: Rust, \texttt{fuser} (FUSE bindings), \texttt{libfuse3}/\texttt{macFUSE}. Verified against a
real local FUSE mount on Linux; all benchmark and bug-reproduction numbers in this report were
measured against real disk images, not simulated.

\end{document}
